\paragraph{同余类残差、审计容量硬界、可见算术计数与零温地基层}
前述 Dirichlet 扭曲谱、审计偶数窗上的 dyadic 容量、有限分辨率可见算术以及输出位势的零温极限，还可继续压缩为下列五条直接后果。周期层的同余类平方根障碍已由定理 \ref{thm:xi-dirichlet-mode-forces-congruence-residual} 给出；下条表明，经 Witt/M\"obius 投影之后，这一障碍在 primitive 层并不会消失；最后一条则把地基层熵、端点极点尺度与零温矩阵的谱半径压缩为同一条平方律。

\begin{theorem}[Witt/M\"obius 投影下的 primitive 同余类平方根障碍]\label{thm:xi-primitive-congruence-residual-nonrh}
沿用命题 \ref{prop:real-input-40-output-dirichlet-nonrh} 与定理 \ref{thm:xi-dirichlet-mode-forces-congruence-residual} 的记号。对每个模数 \(m\ge 2\) 与剩余类 \(a\in\ZZ/m\ZZ\)，定义 primitive 层的同余计数
\[
P_n(a;m):=\#\{\mathcal O:\ \mathcal O\text{ 为长度 }n\text{ 的 primitive 轨道},\ E(\mathcal O)\equiv a\pmod m\},
\qquad
P_n:=\sum_{a\in\ZZ/m\ZZ}P_n(a;m).
\]
若
\[
\theta_m:=\max_{1\le j\le m-1}\frac{\log\rho_{m,j}}{\log\lambda}>\frac12,
\qquad
\lambda=\rho(M(1))=\varphi^2,
\]
则存在某个剩余类 \(a_\ast\in\ZZ/m\ZZ\) 使
\[
\limsup_{n\to\infty}
\frac{1}{n\log\lambda}
\log\Bigl|P_n(a_\ast;m)-\frac{P_n}{m}\Bigr|
\ge
\theta_m
>
\frac12.
\]
因此，在该模数下，primitive 层的同余类分布同样不可能满足任何 RH 形的平方根指数误差界。
\end{theorem}
\begin{proof}
取 \(j_\ast\in\{1,\dots,m-1\}\) 使
\[
\rho_{m,j_\ast}=\lambda^{\theta_m}=:\rho>\lambda^{1/2}.
\]
定义 primitive Fourier 系数
\[
\widehat P_n(j):=\sum_{a\in\ZZ/m\ZZ}P_n(a;m)\omega_m^{ja}.
\]
由命题 \ref{prop:real-input-40-output-dirichlet-nonrh} 中的 Witt/M\"obius 反演，
\[
\widehat P_n(j_\ast)
=
p_n(\omega_m^{j_\ast})
=
\frac{1}{n}\sum_{d\mid n}\mu(d)\,a_{n/d}(\omega_m^{j_\ast d}).
\]
另一方面，定理 \ref{thm:xi-dirichlet-mode-forces-congruence-residual} 的证明已给出
\[
\limsup_{n\to\infty}|a_n(\omega_m^{j_\ast})|^{1/n}=\rho,
\]
故存在子序列 \(n_r\to\infty\) 满足
\[
|a_{n_r}(\omega_m^{j_\ast})|=\rho^{\,n_r+o(n_r)}.
\]

现估计 Möbius 反演中的其余约数项。对任意 \(d\ge 2\)，由逐项取绝对值与 Perron 单调性，
\[
\rho\!\bigl(M(\omega_m^{j_\ast d})\bigr)\le \rho(M(1))=\lambda,
\]
从而
\[
a_{n_r/d}(\omega_m^{j_\ast d})=O(\lambda^{n_r/d})=O(\lambda^{n_r/2})=o(\rho^{n_r}),
\]
因为 \(\rho>\lambda^{1/2}\)。又约数函数满足 \(\tau(n_r)=e^{o(n_r)}\)，故
\[
\sum_{\substack{d\mid n_r\\ d\ge 2}}\mu(d)\,a_{n_r/d}(\omega_m^{j_\ast d})=o(\rho^{n_r}).
\]
于是沿该子序列有
\[
\widehat P_{n_r}(j_\ast)
=
\frac{1}{n_r}a_{n_r}(\omega_m^{j_\ast})+o\!\left(\frac{\rho^{n_r}}{n_r}\right),
\]
进而
\[
\limsup_{n\to\infty}|\widehat P_n(j_\ast)|^{1/n}\ge \rho.
\]

记 primitive 层的同余残差为
\[
b_n^{\mathrm{prim}}(a;m):=P_n(a;m)-\frac{P_n}{m}.
\]
则对 \(j\neq 0\) 有
\[
\sum_{a\in\ZZ/m\ZZ} b_n^{\mathrm{prim}}(a;m)\,\omega_m^{ja}
=
\widehat P_n(j),
\]
因为 \(\sum_a\omega_m^{ja}=0\)。故
\[
|\widehat P_n(j_\ast)|
\le
\sum_{a\in\ZZ/m\ZZ}|b_n^{\mathrm{prim}}(a;m)|
\le
m\max_{a\in\ZZ/m\ZZ}|b_n^{\mathrm{prim}}(a;m)|.
\]
取 \(n\to\infty\) 的上极限并开 \(n\) 次方，得到
\[
\limsup_{n\to\infty}
\left(
\max_{a\in\ZZ/m\ZZ}\Bigl|P_n(a;m)-\frac{P_n}{m}\Bigr|
\right)^{1/n}
\ge
\rho=\lambda^{\theta_m}.
\]
由于剩余类只有有限个，必存在某个 \(a_\ast\in\ZZ/m\ZZ\) 使相应单类残差的上极限至少为 \(\lambda^{\theta_m}\)。两边取对数并除以 \(n\log\lambda\)，即得所示结论。证毕。
\end{proof}

\leanverified{paper\_xi\_foldbin\_audited\_window\_oracle\_hard\_cap}
\begin{theorem}[审计偶数窗上最小退化扇区强制预言机成功率硬上界]\label{thm:xi-foldbin-audited-window-oracle-hard-cap}
沿用定理 \ref{thm:xi-foldbin-oracle-capacity-central-spectral-truncation} 的记号，并定义
\[
\mathrm{Succ}_m^{\mathrm{bin}}(B):=\frac{C_m(B)}{2^m},
\qquad
(x)_+:=\max\{x,0\}.
\]
若 \(m\in\{6,8,10,12\}\)，则对任意整数 \(B\ge 0\) 都有
\[
\boxed{\;
\mathrm{Succ}_m^{\mathrm{bin}}(B)
\le
1-\frac{F_m}{2^m}\bigl(F_{m/2}-2^B\bigr)_+.\;}
\]
特别地，当 \(B<\log_2 F_{m/2}\) 时，
\[
1-\mathrm{Succ}_m^{\mathrm{bin}}(B)
\ge
\frac{F_m}{2^m}\bigl(F_{m/2}-2^B\bigr)
>
0.
\]
\end{theorem}
\begin{proof}
由定理 \ref{thm:fold-bin-min-degeneracy-fib-audited}，当 \(m\in\{6,8,10,12\}\) 时，恰有 \(F_m\) 条纤维达到最小大小
\[
d_{\min}(m)=F_{m/2}.
\]
在这 \(F_m\) 条纤维上，容量公式中的每一项至多贡献 \(2^B\) 个可恢复微态；因此它们总共至多贡献
\[
F_m\min(F_{m/2},2^B).
\]
其余纤维即使全部完全恢复，也至多贡献
\[
2^m-F_mF_{m/2}
\]
个微态。于是
\[
\begin{aligned}
C_m(B)
&\le
\bigl(2^m-F_mF_{m/2}\bigr)+F_m\min(F_{m/2},2^B)\\
&=
2^m-F_m\bigl(F_{m/2}-2^B\bigr)_+.
\end{aligned}
\]
两边除以 \(2^m\) 即得所示成功率上界。证毕。
\end{proof}

\begin{theorem}[有限分辨率可见算术的单位元与幂等元计数闭式]\label{thm:xi-visible-arithmetic-unit-idempotent-count}
\leanverified{paper\_xi\_visible\_arithmetic\_unit\_idempotent\_count}
沿用定理 \ref{thm:fold-congruence-mod-semirings} 的有限分辨率可见算术记号，记
\[
R_m^{\mathrm{vis}}:=\Omega_m/\!\equiv_m
\cong
\ZZ/(F_{m+2}\ZZ).
\]
则其乘法结构满足
\[
\boxed{\;
\#\bigl(R_m^{\mathrm{vis}}\bigr)^\times
=
\phi_{\mathrm{Euler}}(F_{m+2}),\;}
\]
\[
\boxed{\;
\#\{e\in R_m^{\mathrm{vis}}:\ e^2=e\}
=
2^{\omega(F_{m+2})},\;}
\]
其中 \(\phi_{\mathrm{Euler}}\) 为 Euler totient，\(\omega(N)\) 表示 \(N\) 的不同素因子个数。
\end{theorem}
\begin{proof}
由定理 \ref{thm:fold-congruence-mod-semirings}，
\[
R_m^{\mathrm{vis}}\cong \ZZ/(F_{m+2}\ZZ)
\]
作为半环同构成立。故其乘法单位元类与模 \(F_{m+2}\) 的单位元一一对应，数量正是
\[
\phi_{\mathrm{Euler}}(F_{m+2}).
\]

再把 Fibonacci 模数分解为
\[
F_{m+2}=\prod_{i=1}^{r}p_i^{\nu_i},
\qquad
r=\omega(F_{m+2}).
\]
中国剩余定理给出
\[
\ZZ/(F_{m+2}\ZZ)\cong \prod_{i=1}^{r}\ZZ/(p_i^{\nu_i}\ZZ).
\]
在每个局部因子 \(\ZZ/(p_i^{\nu_i}\ZZ)\) 中，幂等方程 \(e^2=e\) 只有 \(e=0,1\) 两个解；因此全局幂等元可在每个坐标独立选取 \(0\) 或 \(1\)，总数恰为 \(2^r\)。这与定理 \ref{thm:xi-fold-congruence-semiring-singularity-nilpotent-profile} 中给出的幂等元计数完全一致。证毕。
\end{proof}

\begin{theorem}[Fibonacci 同余拓扑与 profinite 拓扑完全一致]\label{thm:xi-visible-arithmetic-fibonacci-adic-profinite-coincidence}
\leanverified{paper\_xi\_visible\_arithmetic\_fibonacci\_adic\_profinite\_coincidence}
沿用定理 \ref{thm:fold-congruence-mod-semirings} 的记号。令
\[
\mathcal N_F:=\{F_{m+2}\ZZ:\ m\ge 1\}
\]
并令 \(\tau_F\) 为使 \(\mathcal N_F\) 中每个理想都成为 \(0\) 邻域的最弱线性拓扑。则 \(\tau_F\) 与 \(\ZZ\) 上的 profinite 拓扑 \(\tau_{\mathrm{pro}}\) 完全一致。因而 \(\tau_F\)-完备化满足规范拓扑环同构
\[
\boxed{\;
\widehat{\ZZ}_F\cong \widehat{\ZZ}.\;}
\]
\end{theorem}
\begin{proof}
先证 \(\tau_F\preceq\tau_{\mathrm{pro}}\)。对每个 \(m\ge 1\)，理想
\[
F_{m+2}\ZZ
\]
本身就是 profinite 邻域基 \(\{N\ZZ:\ N\ge 1\}\) 中的一项，故由 \(\mathcal N_F\) 生成的拓扑不强于 profinite 拓扑。

再证反向包含。任取 \(N\ge 1\)。考虑 Fibonacci 递推在有限集
\[
(\ZZ/N\ZZ)^2
\]
上诱导的状态演化
\[
(F_k,F_{k+1})\longmapsto (F_{k+1},F_{k+2}).
\]
由于状态空间有限，轨道必最终周期；而初态 \((F_0,F_1)=(0,1)\) 在周期中必再现。故存在某个 \(L=L(N)\ge 1\) 使
\[
(F_L,F_{L+1})\equiv (0,1)\pmod N.
\]
于是 \(N\mid F_L\)。若 \(N>1\)，则必有 \(L\ge 3\)，从而可写 \(L=m+2\)；对 \(N=1\) 则结论平凡。因而对任意 \(N\ge 1\) 都存在 \(m\ge 1\) 使
\[
F_{m+2}\ZZ\subseteq N\ZZ.
\]
这说明每个 profinite 邻域都包含某个 Fibonacci 邻域，故 \(\tau_{\mathrm{pro}}\preceq\tau_F\)。

综上 \(\tau_F=\tau_{\mathrm{pro}}\)。两种线性拓扑既然一致，其 Hausdorff 完备化作为同一连续有限商系统的逆极限必自然同构，即
\[
\widehat{\ZZ}_F\cong \widehat{\ZZ}.
\]
证毕。
\end{proof}

\begin{corollary}[可见同余半环族对全部有限模算术的共终性]\label{cor:xi-visible-arithmetic-fibonacci-cofinal-quotients}
\leanverified{paper\_xi\_visible\_arithmetic\_fibonacci\_cofinal\_quotients}
沿用定理 \ref{thm:xi-visible-arithmetic-fibonacci-adic-profinite-coincidence} 与定理 \ref{thm:fold-congruence-mod-semirings} 的记号。对任意正整数 \(N\)，存在某个 \(m\ge 1\) 与满射半环同态
\[
\Psi_{m,N}:R_m^{\mathrm{vis}}=\Omega_m/\!\equiv_m\longtwoheadrightarrow \ZZ/(N\ZZ).
\]
等价地，有限半环族
\[
\{\,\Omega_m/\!\equiv_m\,\}_{m\ge 1}
\]
在所有整数商环 \(\ZZ/(N)\) 的意义下是共终的：每个 \(\ZZ/(N)\) 都是某个 \(\Omega_m/\!\equiv_m\) 的商。
\end{corollary}
\begin{proof}
由定理 \ref{thm:xi-visible-arithmetic-fibonacci-adic-profinite-coincidence} 的证明，对任意 \(N\ge 1\) 可取 \(m\ge 1\) 使
\[
N\mid F_{m+2}.
\]
于是标准降模映射给出满射环同态
\[
\rho_{m,N}:\ZZ/(F_{m+2}\ZZ)\longtwoheadrightarrow \ZZ/(N\ZZ).
\]
再由定理 \ref{thm:fold-congruence-mod-semirings}，
\[
R_m^{\mathrm{vis}}=\Omega_m/\!\equiv_m\cong \ZZ/(F_{m+2}\ZZ),
\]
将 \(\rho_{m,N}\) 经此同构拉回，即得所需的满射半环同态 \(\Psi_{m,N}\)。证毕。
\end{proof}

\begin{theorem}[输出位势零温极限的正熵四态地基层与饱和密度上界]\label{thm:xi-real-input-40-output-freezing-positive-entropy-halffill}
\leanverified{paper\_xi\_real\_input\_40\_output\_freezing\_positive\_entropy\_halffill}
沿用命题 \ref{prop:real-input-40-pressure-freezing}、推论 \ref{cor:real-input-40-ground-entropy} 与定理 \ref{thm:killo-real-input-40-positive-entropy-freezing} 的记号。记
\[
P(\theta):=\log\lambda(e^\theta).
\]
则端点斜率满足
\[
\lim_{\theta\to-\infty}P'(\theta)=0,
\qquad
\lim_{\theta\to+\infty}P'(\theta)=\frac12.
\]
并且零温缩放极限对应于四态一步 SFT \(\Sigma_\infty\)，其转移矩阵为
\[
B_\infty=
\begin{pmatrix}
0&1&1&0\\
0&0&1&0\\
0&0&3&1\\
1&0&0&3
\end{pmatrix},
\]
其地基层熵满足
\[
h_{\mathrm{ground}}
=
\log c
=
\frac12\log\rho(B_\infty)
>
0.
\]
因此，极端正偏置下的冻结端面不可能塌缩为有限个周期轨道的并；同时，可见输出 \(1\) 的典型密度始终被压在区间 \([0,\tfrac12]\) 内。
\end{theorem}
\begin{proof}
命题 \ref{prop:real-input-40-pressure-freezing} 已给出
\[
\lim_{\theta\to-\infty}P'(\theta)=0,
\qquad
\lim_{\theta\to+\infty}P'(\theta)=\frac12,
\]
故极端偏置下的典型输出 \(1\) 密度只能落在 \([0,\tfrac12]\) 内。

另一方面，推论 \ref{cor:real-input-40-ground-entropy} 与定理 \ref{thm:killo-real-input-40-positive-entropy-freezing} 说明：饱和端点对应的零温地基并非有限个最大平均能量周期轨道，而是由上式 \(B_\infty\) 所定义的四态 SFT 支配，且其地基层熵满足
\[
h_{\mathrm{ground}}=\log c=\frac12\log\rho(B_\infty)>0.
\]
若冻结端面仅由有限个周期轨道支撑，则其拓扑熵必为 \(0\)，与 \(h_{\mathrm{ground}}>0\) 矛盾。故极端正偏置极限不可能塌缩为有限周期轨道并。证毕。
\end{proof}

\begin{corollary}[正负偏置零温极限分别冻结于 half-filling 与空相]\label{cor:xi-real-input-40-output-freezing-half-filling-empty-phase}
沿用定理 \ref{thm:xi-real-input-40-output-freezing-positive-entropy-halffill} 的记号。记 \(\mu_\theta\) 为输出位势 \(P(\theta)\) 所对应的平衡态族，并记
\[
\mu_\infty^{+}:=\lim_{\theta\to+\infty}\mu_\theta,
\qquad
\mu_\infty^{-}:=\lim_{\theta\to-\infty}\mu_\theta
\]
为两端零温极限测度。则对观测量 \(\mathbf 1_{\mathrm{out}=1}\) 有
\[
\boxed{\;
\int \mathbf 1_{\mathrm{out}=1}\,\dd\mu_\infty^{+}
=
\frac12,\qquad
\int \mathbf 1_{\mathrm{out}=1}\,\dd\mu_\infty^{-}
=
0.\;}
\]
换言之，正偏置零温极限并不饱和到全占据，而是冻结于严格的 half-filling 地基态；负偏置零温极限则冻结于空相。
\end{corollary}
\begin{proof}
对一参数平衡态族，压力导数给出相应观测量的平衡态期望，即
\[
P'(\theta)=\int \mathbf 1_{\mathrm{out}=1}\,\dd\mu_\theta.
\]
由定理 \ref{thm:xi-real-input-40-output-freezing-positive-entropy-halffill}，
\[
\lim_{\theta\to-\infty}P'(\theta)=0,
\qquad
\lim_{\theta\to+\infty}P'(\theta)=\frac12.
\]
令 \(\theta\to\pm\infty\) 即得
\[
\int \mathbf 1_{\mathrm{out}=1}\,\dd\mu_\infty^{-}=0,
\qquad
\int \mathbf 1_{\mathrm{out}=1}\,\dd\mu_\infty^{+}=\frac12.
\]
又同一定理已把 \(\mu_\infty^{+}\) 识别为由四态矩阵 \(B_\infty\) 的正熵地基层所支配的零温极限，故其饱和值严格停在 \(1/2\) 而非 \(1\)。证毕。
\end{proof}
\leanverified{paper\_xi\_real\_input\_40\_output\_freezing\_half\_filling\_empty\_phase}

\begin{theorem}[地基层熵、端点极点尺度与零温矩阵谱半径的平方律]\label{thm:xi-real-input-40-ground-entropy-endpoint-square-law}
沿用定理 \ref{thm:xi-real-input-40-output-freezing-positive-entropy-halffill} 的记号，并记
\[
z_\star(u):=\lambda(u)^{-1},
\qquad
b:=\rho(B_\infty)^{-1}.
\]
则有精确恒等式
\[
\boxed{\;
h_{\mathrm{ground}}
=
\log c
=
-\frac12\log b
=
\frac12\log\rho(B_\infty).\;}
\]
并且当 \(u\to+\infty\) 时，
\[
\boxed{\;
u\,z_\star(u)^2\to b=e^{-2h_{\mathrm{ground}}},
\qquad
z_\star(u)\sim e^{-h_{\mathrm{ground}}}u^{-1/2},
\qquad
\lambda(u)\sim e^{h_{\mathrm{ground}}}u^{1/2}.\;}
\]
\end{theorem}
\begin{proof}
由定理 \ref{thm:xi-real-input-40-output-freezing-positive-entropy-halffill}，
\[
h_{\mathrm{ground}}=\log c=\frac12\log\rho(B_\infty).
\]
故
\[
b=\rho(B_\infty)^{-1}=c^{-2}=e^{-2h_{\mathrm{ground}}},
\]
即得第一条盒装恒等式。

另一方面，命题 \ref{prop:real-input-40-pressure-freezing} 给出
\[
\lambda(u)=c\,u^{1/2}(1+o(1))
\qquad(u\to+\infty).
\]
取倒数便得
\[
z_\star(u)=\lambda(u)^{-1}=c^{-1}u^{-1/2}(1+o(1))
=
e^{-h_{\mathrm{ground}}}u^{-1/2}(1+o(1)).
\]
于是
\[
u\,z_\star(u)^2\to c^{-2}=b=e^{-2h_{\mathrm{ground}}},
\]
再将上式取倒数即得
\[
\lambda(u)\sim e^{h_{\mathrm{ground}}}u^{1/2}.
\]
证毕。
\end{proof}

\endinput
